% 候选集 S={(0,0),(1,0),(0,1)} 的凸包（三角形）上最大化 x1+x2：最优值1，最优面是连接(1,0)与(0,1)的边，含分数点(1/2,1/2)。
\begin{tikzpicture}[x=26mm,y=26mm,
  note/.style={font=\FigureNoteFont\fontsize{7.5}{9}\selectfont,text=FigureMuted}]
  \fill[FigureBlueFill] (0,0)--(1,0)--(0,1)--cycle;
  \begin{scope}\clip (-0.05,-0.05) rectangle (1.35,1.3);
    \foreach \k in {0.5,1.5}{\draw[FigureGreen,line width=.6pt,dashed] (0,\k)--(\k,0);}
  \end{scope}
  \draw[FigureRed,line width=2pt] (1,0)--(0,1);
  \draw[->,FigureStroke,line width=.6pt] (-0.08,0)--(1.4,0) node[right]{$x_1$};
  \draw[->,FigureStroke,line width=.6pt] (0,-0.08)--(0,1.3) node[above]{$x_2$};
  \foreach \p in {(0,0),(1,0),(0,1)}{\fill[FigureBlue] \p circle (2pt);}
  \draw[FigureRed,fill=white,line width=.8pt] (0.5,0.5) circle (1.8pt);
  \node[note,anchor=north east] at (0.47,0.47) {$(1/2,1/2)$};
  \node[note,anchor=north] at (1,-0.03) {$(1,0)$};
  \node[note,anchor=east] at (-0.03,1) {$(0,1)$};
  \node[note,anchor=north east] at (-0.01,-0.01) {$(0,0)$};
  \node[note,anchor=west,text=FigureRed!80!FigureInk] at (0.66,0.5) {最优面};
  \node[note,anchor=west,text=FigureGreen!75!FigureInk] at (1.1,0.42) {$x_1+x_2=3/2$};
  \node[note,text=FigureBlue!80!FigureInk] at (0.3,0.25) {$\operatorname{conv}(S)$};
\end{tikzpicture}
